% Generated by Claude Opus 5 (Anthropic) — compressed §2 and §3 for the 4-page limit.
% §2: 1278 -> ~400 words.  §3: 999 -> ~545 words.

\section{Methodology}
\label{sec:method}

FISH builds a model of a deployed ROS~2 application as the ROS~2 and CUDA
runtimes execute it. Four kinds of observation, on one clock, suffice for this.
\emph{Object identity}: every \texttt{rcl} handle with its creation and
destruction, so that callbacks are attributed to the right node and executor
even when handles are reused. \emph{Dispatch}: the start and end of every
callback execution, whatever its trigger. \emph{Data flow}: which callback
publishes to which topic and which callback consumes it, including deliveries
that never leave the process. \emph{GPU work}: every runtime API call and the
kernels, copies, memsets and synchronisations it produces.

\noindent
\textbf{Instrumentation:}
FISH builds on the two-phase, handle-matched tracepoints of
\texttt{ros2\_tracing}~\cite{bedard-2022-ros2tracing} and enables 42 events: 14
stock and 28 added. The additions fall into six groups. Six finalize events
close the object lifetimes that the stock \texttt{*\_init} events open, and a
seventh binds a timer to its node and period atomically; together they let a
callback be attributed to the right object even after a handle is reused.
Callback start and end at \texttt{execute\_waitable} and
\texttt{execute\_client}, with two client round-trip and four action-server
events, complete the five executable kinds of the \texttt{rclcpp} executor.
Three link events, emitted at the publish, receive and client-request sites,
turn a publish into an edge from the callback that issued it. The intra-process
subscription waitable and two NITROS links cover the deliveries that never reach
\texttt{rcl}. Five registration events give Python nodes the same callback chain
as \texttt{rclcpp}, except for waitables. Four introspection events record the
executor's type and thread count and each callback's group membership. GPU work
is recorded by CUPTI through \texttt{nsys} with a minimal profile and aligned to
the LTTng clock at trace ingest. Correlation ids then bind each GPU record to
its API call, and each API call to the callback window that issued it.

\noindent
\textbf{Acquisition:}
No application source modification is required: the instrumented ROS~2 libraries
are installed on the host or in a containerised environment, and the
\texttt{ros2} entry point is replaced by the FISH entry point, which wraps the
original command. Because CUPTI observes only processes started under the
profiler, the entry point first walks the launch description in describe-only
mode and launches the containers that will host GPU components directly under
\texttt{nsys}. This must happen at spawn time: launch loads composable nodes into
a container only once, so a container restarted later would come up empty.
Standalone GPU processes are handled by the FISH daemon, which polls
\texttt{/proc} and restarts a process under the profiler when it detects a CUDA
context. A system-stable signal then gates the session: the daemon declares the
system up when the node list has been unchanged over five polls and all
components are loaded, and only then starts the input: a rosbag replay, or for a
live system a stimulus schedule, a list of timed commands issued at fixed
offsets from the system-stable instant. Phase boundaries are derived from the
trace itself, so that comparisons use the steady state regardless of how the
input arrived: initialisation ends when every periodic timer has fired once,
warm-up when every recurring callback has completed its first execution.

\noindent
\textbf{From trace to model:}
Data ingest places both traces in one time base. The model generator builds the
hierarchy process--executor--node--entity--callback, attaches the interaction
edges of Section~\ref{sec:model}, and for every callback that submits GPU work
derives a typed DAG per invocation.

\section{Model Definition}
\label{sec:model}

The FISH model follows the abstraction levels of Section~\ref{sec:intro} and
adds the two ROS~2 does not name: the host that carries the processes, and the
entity between a node and its callbacks (Fig.~\ref{fig:model}a). An entity is
what the executor waits on, a callback what it then runs; hence the five entity
types. A publisher is waited on by nobody, and a callback of any kind may
publish, so it is an aspect of the entity whose callback publishes through it.
Interaction is not containment: interactions are edges between callbacks, and
read at a higher layer an edge aggregates those beneath it, so one graph answers
questions at the granularity they are asked (Fig.~\ref{fig:model}b). Namespaces
group vertices for presentation only; an edge crossing a group is lifted to its
boundary as between layers.

\begin{align*}
  v &= (t_v,\ \mathit{id}_v,\ A_v,\ Z_v,\ \mathit{ns}_v), &
  t_v &\in \langle H,\mathit{EX},N,E,F\rangle,\\[2pt]
  e &= (u,\ v,\ \tau,\ \mu), & u,v &\in V_4,\\[2pt]
  E_\ell &= \bigl\{\bigl(\pi(u),\pi(v),\{(\tau,\mu)\}\bigr):\pi(u)\neq\pi(v)\bigr\},
  & \pi &= p^{\,4-\ell}.
\end{align*}

with $t_v$ the layer type, $A_v$ the attributes it carries
(Fig.~\ref{fig:model}a), $Z_v$ the children, $\mathit{ns}_v$ the namespace,
$\tau$ the interaction type, $\mu$ the channel and the per-occurrence
measurements, and $p$ the parent map.

\noindent
\textbf{Edges:}
Two sources produce them. Initialisation declares the dispatch wiring, and with
it the edges that release a callback: $\mathit{msg}$ where a publisher and a
subscription meet on a topic, and a $\mathit{dep}$ pair, request and response,
where a client and a service meet on a service name. Execution completes these
twice over --- it resolves the source end to a callback, which registration
cannot do, and it records the delivery path and, for $\mathit{msg}$, the hop
latency of every occurrence. The other two edges exist only in execution.
$\mathit{state}$ is a sample-and-hold read: the destination takes the latest
value the source left in the node rather than a message delivered to it, so it
carries data age instead of hop latency, and it is $\mathit{polled}$ when the
read is seen as an explicit take and $\mathit{inferred}$ otherwise.
$\mathit{join}$ is a synchronisation constraint rather than a data edge: a
node's output is produced by whichever of several inputs completes the set, and
$\mathit{join}$ ties the set to that completer. Figure~\ref{fig:model}c gives
the trace pattern each is decided from. An untraced endpoint or an out-of-ROS
thread becomes an $\mathit{external}$ vertex.

\noindent
\textbf{GPU DAG:}
A callback that submits GPU work owns one typed DAG per invocation: CPU
segments, compute-engine kernels and copy-engine transfers, ordered by launch,
host order and stream FIFO. FISH groups invocations by signature into shapes and
folds them into a conditional graph (Fig.~\ref{fig:model}d) --- the backbone
common to all executions plus the observed branches with their shares --- which
is the input-dependent structure a GPU-aware analysis must bound.

\noindent
\textbf{Correctness:}
A map $\chi$ sends every vertex-creating ROS~2 API call to the records it must
emit and the vertices it must create; a node constructor, for instance, yields
$1N+7E+7F$, since it also creates six parameter services and the time-source
subscription. Correctness is then $\chi_r(A)\subseteq R$ and $\chi_v(A)\subseteq
V$ over the executed calls $A$: every expected record is in the trace, and every
expected vertex in the graph.
